\section{Related work}
\label{sec:related}

\subsection{Truncation-degree selection}

\ph{one paragraph} Attenuation-only thresholds, degree-variance criteria in the
Meissl / Lowes--Mauersberger form \citep{RummelVanGelderen1995,Lowes1966}, and
Kaula-type spectral rules \citep{kaula1966theory}. Radius-driven selection
during propagation \citep{atallah2022radial}. These methods commonly select on
a local criterion --- force error, radius, spectral tail, or a state or
uncertainty proxy --- evaluated pointwise. The distinction drawn here is not
that such criteria are unknown to the literature but that none of them
allocates the \emph{degree history} jointly over a finite horizon under one
budget while retaining the signed cross-epoch coupling of
Eq.~\eqref{eq:objective}; that is the gap Section~\ref{sec:problem}
formalizes.

\subsection{Budgeted allocation}
\label{sec:related-allocation}

The separable case is classical. Allocating a scalar budget over independent
items through a single multiplier is Everett's generalized Lagrange multiplier
method \citep{everett1963generalized}, and the discrete form used for a set of
quantizers is the standard treatment in rate--distortion-style bit allocation
\citep{shoham1988efficient}, which also records that over a discrete set the
multiplier sweep recovers only supported efficient solutions, so a duality gap
may remain. The previous paper used exactly that construction at force level.

What changes here is that the objective couples the items, so neither result
applies directly. The coupled form of Eq.~\eqref{eq:allocation} is an integer
quadratic program with a single knapsack constraint, a class with its own
literature \ph{position against the quadratic-knapsack and integer-QP
literature; the citation pass supplies the references}. What that literature
supplies is the expectation of what is achievable --- a good feasible point and
a relaxation bound rather than a certified optimum --- and that is what
Section~\ref{sec:algorithm} delivers. What it does not supply is the
$\max(i,k)$ structure of Eq.~\eqref{eq:Q-structure}, which is specific to the
objective being a quadrature over prefix sums and is what makes a sweep linear
rather than quadratic in the number of epochs. From the separable literature,
what carries over is the language: an \emph{attainable} allocation reported as
such rather than as a bound.

\subsection{Goal-oriented error control}

Weighting a local residual by the sensitivity of a functional of interest is
the dual-weighted-residual idea in adaptive finite elements
\citep{becker2001optimal}, and the distinction it draws --- between a residual
norm and an adjoint-weighted residual --- is the same distinction the previous
paper measured in this setting. Two things separate the present problem from
that literature. The functional here is quadratic in an integral of the
residual rather than linear in it, so the weighting is not a scalar field;
and the control variable is a discrete degree with a quadratic cost, so the
adaptivity decision is a budgeted integer allocation rather than a refinement
flag.

\subsection{Discrete optimal control}

The structure this paper exploits is, read from the continuous side, an
adjoint one, and it is worth saying so rather than leaving a reader to notice
it. Treating the degree as a discrete control and the accumulated displacement
as a state, the allocation is a Mayer--Bolza problem with a mixed integer
control; Pontryagin's principle \citep{pontryagin1962mathematical} then
reduces the choice at each instant to minimizing a Hamiltonian against a
costate, and the cross term $\mathbf c_i$ of Eq.~\eqref{eq:cross-term} is
proportional to that costate transported to $t_0$ --- one half of
$\nabla_{\mathbf u_i}J=2\mathbf c_i$ under the normalization used here, the
factor depending on the adjoint convention. The forward--backward iteration of
Section~\ref{sec:controller-ifbda} is the corresponding two-point boundary
value problem solved by successive sweeps
\citep[Ch.~2]{bryson1975applied}, and the same reading places the method in
the mixed-integer optimal control literature, where relaxation followed by
sum-up rounding is the standard device \citep{sager2012integer}.

Naming the connection costs nothing and clarifies what is and is not new here.
The adjoint is not: it is a century old. What is specific is the combination
--- a spherical-harmonic truncation degree as the control, an objective that
is quadratic in a \emph{signed} integral so that the coupling is genuinely
cross-epoch, a discretization whose coupling matrix collapses onto
$\max(i,k)$ and therefore admits a certified bound rather than only a
stationarity condition, and a deployable approximation whose limiting
difficulty is informational rather than computational.

\subsection{Sensitivity in astrodynamics}

The state-transition matrix and the variational equations that generate it are
standard \citep[Ch.~4]{tapley2004statistical}. Mapping geopotential
coefficients to radial, transverse and normal position error analytically about
a Kepler orbit is classical \citep{rosborough1987radial}. \ph{one or two
sentences} on what is different here: the mapping is used forward, as an
allocation weight, and it is driven by a truncation defect rather than by an
estimation residual or a single coefficient.

\subsection{Multifidelity models}

\ph{short paragraph} \citep{peherstorfer2018multifidelity,
jones2019multifidelity}. The degree schedule is a fidelity schedule, and the
question of where to spend fidelity under a budget is shared; what is specific
here is that the low-fidelity model's error has an exploitable sign structure
along the trajectory.
